Large language models produce text with a distinctive style---formal, hedging, verbose---that readers readily identify as ``sounding like AI.''
Recent work in mechanistic interpretability has shown that high-level concepts such as truthfulness and refusal are linearly represented in the residual stream of transformer models, but whether stylistic qualities like \aistyle follow the same geometric pattern remains unexplored.
We investigate this question using contrastive activation analysis on paired human- and ChatGPT-written text from the \hcthree dataset.
Linear probes trained on residual stream activations achieve 93--96\% accuracy at distinguishing human from AI text across two architecturally different models (\pythia and \qwen), while the one-dimensional difference-in-means direction alone achieves 69--81\% accuracy, significantly above chance.
Cross-domain generalization averages 69--72\%, and causal steering along the identified direction produces monotonic shifts in activation projections.
These results provide evidence that \aistyle is a linearly decodable feature of the residual stream, though the gap between probe and single-direction accuracy suggests it occupies a multi-dimensional subspace rather than a single clean direction.
Our findings connect mechanistic interpretability to AI-generated text detection and suggest that alignment training produces structured, geometrically simple stylistic signatures.
